\par\noindent\textbf{(1)} 线性化方程 \(\partial_t\boldsymbol u+2\boldsymbol\Omega\times\boldsymbol u=-\nabla p/\rho\)，\(\nabla\cdot\boldsymbol u=0\)。对平面波取 \(\boldsymbol u=\boldsymbol Ue^{i(\boldsymbol k\cdot\boldsymbol x-\omega t)}\)，用 \(\boldsymbol k\cdot\boldsymbol U=0\) 消去压力，只留下受Coriolis耦合的横、纵速度。
\par\noindent\textbf{(2)} 代数线性系统行列式给 \(\omega^2=4\Omega^2k_z^2/(k_x^2+k_z^2)\)，所以 \(|\omega|\le2|\Omega|\)。正频率可写 \(\omega=2\Omega k_z/k\)（选定传播支和 \(\Omega>0\)），频率仅由波矢与旋转轴夹角决定而非波长模。
\par\noindent\textbf{(3)} 对 \(\omega=2\Omega k_z/\sqrt{k_x^2+k_z^2}\) 求梯度，得 \(v_{gx}=-2\Omega k_xk_z/k^3\)、\(v_{gz}=2\Omega k_x^2/k^3\)。点积 \(k_xv_{gx}+k_zv_{gz}=0\)，因此理想惯性波群能流垂直相位波矢，和普通各向同性声波截然不同。
\par\noindent\textbf{(4)} 若 \(k_z=0\)，则 \(\omega=0\)，扰动成为静态地转型结构；若 \(k_x=0\)，则 \(|\omega|=2|\Omega|\)，群速度趋零。零波数时方向不定义，不能令 \(k\to0\) 而忽略所取角度；有限边界和黏性会改变惯性模。
